% Figura redibujada a partir de
% "Imagenes/Velocidad densidad de corriente frente a campo electrico.PNG" (Tema 2).
% Ocupa lo mismo que el PNG en el documento: 386.0 pt x 307.1 pt. Las
% posiciones se dan en píxeles del PNG original (641 x 510 px).
% Modelos (T = 300 K):
%   Si, Ge (Caughey-Thomas): v = mu E / (1 + (mu E / v_s)^b)^(1/b)
%   GaAs, electrones:        v = (mu E + v_s (E/E_0)^b) / (1 + (E/E_0)^b)
\documentclass[12pt,border=0pt]{standalone}
\usepackage[spanish,es-tabla]{babel}
\usepackage[utf8]{inputenc}
\usepackage{amsmath}
\usepackage{xcolor}
\usepackage{tikz}
\usepackage{pgfplots}
\usepackage{circuitikz}
\usetikzlibrary{babel, circuits.ee}
\pgfplotsset{compat=1.18}

\definecolor{applegreen}{rgb}{0.55, 0.71, 0.0}
\definecolor{airforceblue}{rgb}{0.36, 0.54, 0.66}
\definecolor{amber}{rgb}{1.0, 0.49, 0.0}
\colorlet{colSi}{blue}
\colorlet{colGe}{black}
\colorlet{colGaAs}{red}

\newcommand{\diff}[3][1]{%
	\IfBlankTF{#1}{\frac{\mathrm{d}{#2}}{\mathrm{d}{#3}}}%
	{\frac{\mathrm{d}^{#1} {#2}}{\mathrm{d}{#3}^{#1}}}%
}%
\newcommand{\pdiff}[3][1]{%
	\IfBlankTF{#1}{\frac{\partial {#2}}{\partial {#3}}}%
	{\frac{\partial^{#1} {#2}}{\partial {#3}^{#1}}}%
}%

\begin{document}
\begin{tikzpicture}[x=0.60225pt, y=-0.60225pt, font=\footnotesize]
    \useasboundingbox (0,0) rectangle (641,510);
    \begin{scope}[x=1pt, y=1pt]
    % Esquina inferior izquierda de los ejes en (78, 443) px
    \begin{axis}[
        at={(46.976pt,-266.80pt)}, anchor=south west, scale only axis,
        width=323.41pt, height=253.55pt,
        xmode=log, ymode=log, xmin=1e2, xmax=1e6, ymin=1e5, ymax=1e8,
        grid=both, major grid style={line width=0.35pt, black!45},
        minor grid style={line width=0.2pt, black!20},
        tick align=inside, tick style={black, line width=0.3pt},
        samples=100,
        declare function={
            ct(\E,\mu,\vs,\b) = \mu*\E/(1 + (\mu*\E/\vs)^\b)^(1/\b);
            gaas(\E) = (8000*\E + 4.5e6*(\E/4500)^2.25)/(1 + (\E/4500)^2.25);
        },
    ]
        % GaAs (electrones)
        \addplot[colGaAs, ultra thick, domain=2:5.3] ({10^x}, {gaas(10^x)});
        % Ge
        \addplot[colGe, ultra thick, domain=2:4.83] ({10^x}, {ct(10^x, 3900, 4.8e6, 2.5)});
        \addplot[colGe, ultra thick, dashed, domain=2:4.83] ({10^x}, {ct(10^x, 2200, 4.8e6, 2.5)});
        % Si
        \addplot[colSi, ultra thick, domain=2:5.12] ({10^x}, {ct(10^x, 1350, 1.02e7, 1.11)});
        \addplot[colSi, ultra thick, dashed, domain=2.3:5.1] ({10^x}, {ct(10^x, 480, 8.4e6, 1.21)});
    \end{axis}
    \end{scope}
    % Etiquetas de los ejes
    \node[rotate=90] at (14,232) {Velocidad de arrastre (cm/s)};
    \node at (345,497) {Campo eléctrico (V/cm)};
    % Rótulos de las curvas
    \node[text=colGaAs, fill=white, inner sep=1pt] at (305,104) {GaAs (electrones)};
    \node[text=colGe, fill=white, inner sep=1pt] (Ge) at (116,218) {Ge};
    \draw[-latex] (Ge.east) to[out=0, in=160] (217,231);
    \draw[-latex] (Ge.east) to[out=-5, in=150] (239,236);
    \node[text=colSi, fill=white, inner sep=1pt] (Si) at (308,322) {Si};
    \draw[-latex] (Si.north) to[out=110, in=-30] (238,268);
    \draw[-latex] (Si.north) to[out=90, in=-80] (300,266);
    % Leyenda
    \fill[white] (385,268) rectangle (530,352);
    \node at (455,285) {$T = 300$ K};
    \draw[ultra thick] (395,312) -- (432,312);
    \node[anchor=west] at (438,312) {Electrones};
    \draw[ultra thick, dashed] (395,338) -- (432,338);
    \node[anchor=west] at (438,338) {Huecos};
\end{tikzpicture}
\end{document}
